\chapter{Transforms, Oscillatory Tails, and Zeta}
\label{ch:transforms-zeta}

The same function can be computed by a transform, an equation, or a finite
sum with endpoint corrections.  Three examples show why the foundation
must keep track of the error appropriate to each operation.  None requires
an ambient space of arbitrary continuous or square-integrable functions.

\section{Laplace transforms with an explicit half-plane}
Suppose a concrete, locally integrable computation $f$ has a rational
exponential bound $|f(t)|\le B E(\alpha t)$ on $t\ge0$.  On a certified
half-plane $\Re s\ge\alpha+\sigma$, with rational $\sigma>0$, define
\[
 \mathcal Lf(s)=\int_0^\infty E(-st)f(t)\,dt.
\]
Its far tail is at most $B E(-\sigma T)/\sigma$.  The middle integral
has the data of the particular $f$.  This constructs the transform.
Multiplication of the integrand by $t^k$ still gives a computable
exponential tail, by the polynomial-exponential calculation in
Chapter~\ref{ch:improper-parameters}.  Expanding $E(-ht)$ with a smaller
half-plane margin gives the parameter-majorant proof of
\[
 (\mathcal Lf)^{(k)}(s)=(-1)^k\int_0^\infty t^kE(-st)f(t)\,dt.
\]
The half-plane is part of the certificate, not an unspecified region of
convergence.

If $f$ is a certified piecewise primitive with finitely many jumps and
its ordinary derivative has the needed tails, finite integration by parts
and the explicit jump contributions give
\[
 \mathcal L(Df)(s)=s\mathcal Lf(s)-f(0+),
\]
where the transform includes interior impulse masses and no impulse at
the initial endpoint.  Extending $f$ by zero to negative times instead
adds $f(0+)\delta_0$ to its full-line derivative.  This distinguishes
the two initial-value conventions rather than assigning a hidden half
weight to the delta at zero.

For an interior impulse $\delta_\tau$, its transform is $E(-s\tau)$.
In particular, ordinary integrations by parts give
\[
 \mathcal L(S)(s)=\frac\pi{s^2+\pi^2},\qquad
 \mathcal L(C)(s)=\frac{s}{s^2+\pi^2}\quad(\Re s>0).
\]
The denominators are separated on any supplied closed box inside that
half-plane.  These formulas compare the oscillator's impulse kernel with
a rational computation; they are not applications of an unproved transform
inversion theorem.

For a constant matrix $A$, the zero-past impulse response is $E(tA)$
for $t>0$.  On $\Re s>\|A\|$, entrywise integration gives
\[
 (sI-A)\mathcal L(E(\cdot A))(s)=I.
\]
The finite geometric matrix series
$s^{-1}\sum_{n\ge0}(A/s)^n$, with its norm tail, independently constructs
$(sI-A)^{-1}$.  Uniqueness of that inverse proves their equivalence.  Thus
the familiar resolvent is a second computation of the fundamental matrix's
Laplace transform.

\section{An oscillatory integral with no absolute-integral bound}
Extend $S(t)/t$ at zero by the series value $\pi$.  For $0<T<U$,
integration by parts using $DS=\pi C$, $DC=-\pi S$ gives
\[
 \left|\int_T^U\frac{S(t)}t\,dt\right|\le\frac{2}{\pi T}.
\]
Indeed the endpoint terms have bounds $1/(\pi T)$ and $1/(\pi U)$,
and the remaining cosine integral is bounded by
$(1/T-1/U)/\pi$.  Therefore bounded quadrature and this cancellation
bound construct $I=\int_0^\infty S(t)/t\,dt$.

To identify it, first use a positive real damping parameter $a$:
\[
 F(a)=\int_0^\infty E(-at)\frac{S(t)}t\,dt.
\]
On $a\ge a_0>0$ the exponential tails justify parameter differentiation,
giving $F'(a)=-\pi/(a^2+\pi^2)$ by the Laplace calculation above.
Also $|F(a)|\le\pi/a$ from the chord bound $|S(t)|\le\pi t$.
Integrating this derivative between positive rational endpoints and sending
the larger one to infinity with that explicit bound yields
\[
 F(a)=\int_a^\infty\frac\pi{u^2+\pi^2}\,du
     =\frac\pi2-A(a/\pi)
\]
when $0<a\le\pi$; scaling and the previously computed angle integral
prove the last equality.

The damping can be removed with a finite estimate.  The function
$E(-at)/t$ is positive decreasing; the same integration-by-parts argument
bounds its oscillatory tail uniformly for all $a\ge0$ by $2/(\pi T)$.
On $[0,T]$, the bound $|1-E(-at)|\le at$ gives
\[
 |F(a)-I|\le aT+\frac4{\pi T}.
\]
First choose a rational $T$, then a sufficiently small positive rational
$a$.  This proves $F(a)\to I$ without an absolute domination on the
infinite interval.  Since $A(a/\pi)\to0$ by its continuity estimate,
\[
 \boxed{\int_0^\infty\frac{S(t)}t\,dt=\frac\pi2.}
\]
For the normalized sinc $S(t)/(\pi t)$, the positive-half-line integral
is $1/2$ and the whole-line integral is $1$.  The endpoint and normalization
are stated explicitly.

\section{A Fourier transform proved by an ODE}
Let $f(x)=E(-\pi x^2)$ and define for real $\xi$
\[
 H(\xi)=\int_{-\infty}^\infty f(x)E(-2\pi i\xi x)\,dx.
\]
The gamma normalization and scaling give $H(0)=1$.  Each parameter
derivative is controlled by a polynomial times the Gaussian.  For complex
increments of size at most $r$, the envelope gains at most $E(2\pi r|x|)$.
Once $|x|\ge4r$, this is still bounded together with the Gaussian by
$E(-\pi x^2/2)$.  The remaining bounded interval is handled separately.
This provides explicit integrable majorants for the parameter series.

Differentiate under the integral by those bounds, and integrate the
identity $f'(x)=-2\pi xf(x)$ by parts on finite symmetric intervals.
Its boundary terms tend to zero with their Gaussian bounds.  The result is
\[
 H'(\xi)=-2\pi\xi H(\xi),\qquad H(0)=1.
\]
The integral-equation uniqueness of the ODE chapter identifies its
computation as
\[
 \boxed{\int_{-\infty}^\infty E(-\pi x^2)E(-2\pi i\xi x)\,dx
                         =E(-\pi\xi^2).}
\]
Multiplying the integrand by $x^k$ gives
$(-2\pi i)^{-k}$ times the $k$th derivative on the right.  Thus each
Gaussian times a specified polynomial has its transform computed by finite
differentiation and polynomial algebra.  This is a genuine Fourier result
without constructing a space of arbitrary $L^2$ functions.

\section{Bernoulli polynomials are finite rational recurrences}
To go beyond a series' original convergence region we shall use finite
summation by parts.  Define polynomials $B_n(x)$ by
\[
 B_0=1,\qquad B_n'=nB_{n-1},\qquad\int_0^1 B_n(x)\,dx=0\quad(n\ge1).
\]
Polynomial integration fixes the constant term uniquely, so every coefficient
is rational.  The first polynomials are
\[
 B_1=x-\tfrac12,\quad B_2=x^2-x+\tfrac16,
 \quad B_4=x^4-2x^3+x^2-\tfrac1{30}.
\]
Induction using the recurrence and the zero mean proves
$B_n(1-x)=(-1)^nB_n(x)$.  Also $B_n(1)=B_n(0)$ for $n\ge2$.
Consequently the odd endpoint values vanish for odd $n>1$.
Put $b_{2k}=B_{2k}(0)$ and $\widetilde B_n(t)=B_n(t-\lfloor t\rfloor)$.
For integration this is simply a separate polynomial on each unit cell;
no discontinuous floor computation at a computed input is needed.

\section{Euler--Maclaurin as a finite identity}
\label{thm:finite-euler-maclaurin}
For integers $N<M$ and a specific function $f$ with certified derivatives
and integration-by-parts identities through order $2p$, one has
\begin{align*}
 \sum_{n=N}^{M-1} f(n)
 &=\int_N^M f(t)\,dt+\frac{f(N)-f(M)}2\\
 &\quad+\sum_{k=1}^p\frac{b_{2k}}{(2k)!}
                \bigl(f^{(2k-1)}(M)-f^{(2k-1)}(N)\bigr)\\
 &\quad-\frac1{(2p)!}\int_N^M\widetilde B_{2p}(t)f^{(2p)}(t)\,dt.
\end{align*}
\begin{proof}
On $[n,n+1]$, integration by parts gives
\[
 \int_n^{n+1}B_1(t-n)f'(t)\,dt
       =\tfrac12(f(n)+f(n+1))-\int_n^{n+1}f(t)\,dt.
\]
Sum these identities.  Repeatedly use
$\widetilde B_j=(\widetilde B_{j+1}/(j+1))'$ inside each cell to integrate
the remaining term by parts.  Matching endpoint values cancel the internal
boundaries for $j+1\ge2$, and the odd endpoint values above eliminate the
odd corrections.  After $2p$ derivatives the displayed remainder is what
remains.  Every step is a finite polynomial-weighted integral identity.
\end{proof}

This is the finite formula underlying the usual
\href{https://dlmf.nist.gov/2.10}{Euler--Maclaurin expansion}.  Its remainder
is part of the computation, not an unspecified asymptotic error.

\section{Zeta beyond the defining sum}
On a certified half-plane $\Re s\ge\sigma>1$, the series
\[
 \zeta(s)=\sum_{n\ge1}E(-sL(n))
\]
has tail at most $N^{1-\sigma}/(\sigma-1)$ after $N$ terms, by finite
comparison with the integral of $t^{-\sigma}$.  This first constructs
zeta only where the series has that bound.

Apply the finite Euler--Maclaurin identity to $f(t)=t^{-s}$, whose
uniform increment and derivative identities follow from the elementary
exponential and logarithm computations.  Here
\[
 f^{(j)}(t)=(-1)^j s^{\overline j}t^{-s-j},\qquad
 s^{\overline j}=s(s+1)\cdots(s+j-1)
\]
is a rising factorial, with $s^{\overline0}=1$.
Sending the upper endpoint to infinity first in $\Re s>1$ gives
\begin{align*}
 \zeta(s)&=\sum_{n=1}^{N-1}n^{-s}
       +\frac{N^{1-s}}{s-1}+\frac12N^{-s}\\
 &\quad+\sum_{k=1}^p\frac{b_{2k}}{(2k)!}
                   s^{\overline{2k-1}}N^{-s-2k+1}
       +R_{N,p}(s),\\
 R_{N,p}(s)&=-\frac{s^{\overline{2p}}}{(2p)!}
              \int_N^\infty\widetilde B_{2p}(t)t^{-s-2p}\,dt.
\end{align*}
Let $C_{2p}$ be the sum of the absolute coefficients of $B_{2p}$, a
rational upper bound on its absolute value in $[0,1]$.  For
$\Re s\ge\sigma>1-2p$, the remainder satisfies
\[
 \boxed{|R_{N,p}(s)|\le
 \frac{|s^{\overline{2p}}|C_{2p}}{(2p)!(\sigma+2p-1)}
                 N^{1-\sigma-2p}.}
\]
A rational upper bound for the finite rising-factorial product is used
for computed input boxes.  The remainder can be omitted to this error,
or computed more accurately by integrating finitely many unit cells and
using the same bound at a larger cutoff.

Thus for a particular $s\ne1$, choose $p$ with $\Re s>1-2p$ and a
certified separation from $s=1$, then increase $N$ until the error is small.
This computes a value even when the original series diverges.
Independence of $N$ follows from the finite summation identity.  Independence
of $p$ follows by two further integrations by parts in the remainder:
the extra correction is precisely the next displayed Bernoulli term,
and the far boundary vanishes in the common half-plane.  This proves
agreement without assuming an analytic continuation theorem.

These formulas are analytic there, with their only possible pole at $1$.
For the remainder, expand $t^{-w}$ near any fixed parameter.  The estimate
$|L(t)|^n/n!\le\rho^{-n}t^\rho$ gives the parameter-majorant proof
whenever $\rho<\Re s+2p-1$.  Hence the constructions also supply local
charts and all requested parameter derivatives.  The residue at $1$ is
$1$, read directly from $N^{1-s}/(s-1)$; its remaining part is computed
by the removable series quotient.

At $s=0$ take $p=1$; its rising factorials vanish and the formula gives
$\zeta(0)=-1/2$.  At $s=-1$ take $p=2$, so the remainder vanishes and
the finite terms give $\zeta(-1)=-1/12$.  This is a value of the proved
continuation computation, not a claim that the divergent sum of positive
integers converges to a negative number.  Standard zeta notation and
continuation formulas are recorded in
\href{https://dlmf.nist.gov/25.2}{DLMF, Section 25.2}.

\section{The common foundation is the finite comparison}
The constructions in this chapter use different kinds of error: an
exponential far tail, oscillatory integration by parts, Gaussian parameter
majorants, and a periodic-polynomial summation remainder.  None can be
replaced by a bare claim that the function is familiar.  Once its bound
has been proved, however, the existing interval, integration, derivative,
and equivalence constructions apply to it without rebuilding the theory.
